% Bagian: sets-functions-relations
% Bab: functions
% Subbagian: isomorphisms

\documentclass[../../../include/open-logic-section]{subfiles}


\begin{document}

\olfileid[id]{sfr}{fun}{iso}
\olsection{Isomorfisme}

\begin{explain}
Suatu \emph{isomorfisme} ialah bijeksi yang mempertahankan struktur
himpunan-himpunan yang dihubungkannya, dengan struktur ditentukan oleh
relasi-relasi yang berlaku di antara !!{element}s himpunan tersebut. Tinjaulah
dua himpunan berikut, $X=\{1,2,3\}$ dan $Y=\{4,5,6\}$. Kedua himpunan ini
memiliki struktur yang diberikan oleh relasi penerus, lebih kecil dari, dan
lebih besar dari. Suatu isomorfisme antara kedua himpunan itu ialah
!!a{bijection} yang mempertahankan struktur-struktur tersebut. Jadi, suatu
fungsi !!a{bijective} $f \colon X \to Y$ merupakan isomorfisme jika $i<j$
jika dan hanya jika $f(i)<f(j)$, $i>j$ jika dan hanya jika $f(i)>f(j)$, dan
$j$ merupakan penerus $i$ jika dan hanya jika $f(j)$ merupakan penerus
$f(i)$.
\end{explain}

\begin{defn}[Isomorfisme]
Misalkan $U$ merupakan pasangan $\langle X, R\rangle$ dan $V$ merupakan
pasangan $\langle Y, S\rangle$, dengan $X$ dan $Y$ himpunan, serta $R$ dan
$S$ masing-masing relasi pada $X$ dan $Y$. Suatu !!{bijection} $f$ dari $X$
ke $Y$ merupakan \emph{isomorfisme} dari $U$ ke $V$ jika dan hanya jika
bijeksi itu mempertahankan struktur relasional, yakni, untuk sembarang $x_{1}$ dan $x_{2}$
dalam $X$, berlaku $\tuple{x_1,x_2} \in R$ jika dan hanya jika
$\tuple{f(x_1),f(x_2)} \in S$.
\end{defn}

\begin{ex}
Tinjaulah dua himpunan $X=\{1,2,3\}$ dan $Y=\{4,5,6\}$, beserta relasi lebih
kecil dari dan lebih besar dari. Fungsi $f\colon X \to Y$ dengan
$f(x) = 7-x$ merupakan isomorfisme antara $\tuple{X,<}$ dan $\tuple{Y,>}$.
\end{ex}

\end{document}
